%============================================================
%  Bd. 44 - Kommutative Ringe mit Eins
%============================================================
\documentclass[main.tex]{subfiles}
\ifSubfilesClassLoaded{\BandSetupStandalone{B44}}{}
\title{Bd. 44 - Kommutative Ringe mit Eins}
\author{Martin Kunze}
\date{}
\setcounter{file}{44}
\hypersetup{pdftitle={Bd. 44 - Kommutative Ringe mit Eins},pdfauthor={Martin Kunze}}
\begin{document}
\providecommand{\BandXXIXProofRefStack}[1]{%
  \ensuremath{\begin{gathered}[t]#1\end{gathered}}}
\title{Bd. 44 - Kommutative Ringe mit Eins}
\author{Martin Kunze}
\date{}
\hypersetup{pageanchor=false}
\maketitle
\hypersetup{pageanchor=true}
\setcounter{tocdepth}{3}
\tableofcontents
\setcounter{file}{44}
\renewcommand{\FormulaBandID}{44}

\chapter{Einleitung}

Dieser Band ergänzt die Ringe mit Eins aus Band~43 um die
Kommutativität der Multiplikation. Die Morphismen erhalten dieselben
Operationen und die Eins wie allgemeine Ringhomomorphismen; zusätzlich
werden Quelle und Ziel als kommutative Ringe vorausgesetzt. Als
grundlegendes Beispiel wird die Kommutativität des Rings der ganzen
Zahlen festgehalten.

\chapter{Kommutative Ringe mit Eins}

\FormulaAxiomDeltaK[Kommutativität der Ringmultiplikation]{%
  \RingAlg{R}{+}{\cdot}{0}{1}\dsep a,b\in R
  \vdash a\cdot b=b\cdot a%
}{CommutativeRingMultiplicationLaw}{%
  \DeltaRow{Mengen}{R\dsep 0\dsep 1\dsep a\dsep b}
  \DeltaRow{Funktionen}{+\dsep\cdot}
}

\FormulaDefDeltaK[Begriff des kommutativen Rings mit Eins]{%
  \CommutativeRingAlg{R}{+}{\cdot}{0}{1}%
}{CommutativeRingDef}{%
  \DeltaRow{Mengen}{R\dsep 0\dsep 1\dsep a\dsep b}
  \DeltaRow{Funktionen}{+\dsep\cdot}
  \DeltaRow{\textbf{Axiome}}{}
  \DeltaRow{Ring mit Eins}{\RingAlg{R}{+}{\cdot}{0}{1}}
    [\FormulaRefAuto{RingDef}[def]]
  \DeltaRow{Kommutative Multiplikation}{%
    a,b\in R\vdash a\cdot b=b\cdot a}
    [\FormulaRefAuto{CommutativeRingMultiplicationLaw}[ax]]
  \DeltaRow{\textbf{Neues Symbol}}{}
  \DeltaPrem{Kommutativer Ring mit Eins}
    {\CommutativeRingAlg{R}{+}{\cdot}{0}{1}}
}

\FormulaAxiomDeltaK[Ringanteil eines kommutativen Rings]{%
  \CommutativeRingAlg{R}{+}{\cdot}{0}{1}
  \vdash\RingAlg{R}{+}{\cdot}{0}{1}%
}{CommutativeRingBaseRingAxiom}{%
  \DeltaRow{Mengen}{R\dsep 0\dsep 1}
  \DeltaRow{Funktionen}{+\dsep\cdot}
}

\FormulaAxiomDeltaK[Multiplikationskommutativität im kommutativen Ring]{%
  \CommutativeRingAlg{R}{+}{\cdot}{0}{1}\dsep a,b\in R
  \vdash a\cdot b=b\cdot a%
}{CommutativeRingMultiplicationAxiom}{%
  \DeltaRow{Mengen}{R\dsep 0\dsep 1\dsep a\dsep b}
  \DeltaRow{Funktionen}{+\dsep\cdot}
}

\chapter{Homomorphismen kommutativer Ringe}

Für kommutative Ringe bleiben die Erhaltungsgleichungen unverändert.
Die Kommutativität gehört zu den beiden beteiligten Ringstrukturen und
ist keine zusätzliche Forderung an die Abbildung.

\FormulaDefDeltaK[Homomorphismen kommutativer Ringe]{%
  \CommutativeRingHom{f}{R,\oplus,\odot,0_R,1_R}{S,\boxplus,\boxtimes,0_S,1_S}%
}{CommutativeRingHomDef}{%
  \DeltaRow{Mengen}{R\dsep S\dsep0_R\dsep1_R\dsep0_S\dsep1_S}
  \DeltaRow{Funktionen}{f\dsep\oplus\dsep\odot\dsep\boxplus\dsep\boxtimes}
  \DeltaRow{\textbf{Axiome}}{}
  \DeltaPrem{Quellstruktur}{\CommutativeRingAlg{R}{\oplus}{\odot}{0_R}{1_R}}
  \DeltaPrem{Zielstruktur}{\CommutativeRingAlg{S}{\boxplus}{\boxtimes}{0_S}{1_S}}
  \DeltaPrem{Homomorphismus}{\RingHom{f}{R,\oplus,\odot,0_R,1_R}{S,\boxplus,\boxtimes,0_S,1_S}}
  \DeltaRow{\textbf{Neues Symbol}}{}
  \DeltaPrem{Homomorphismus}{\CommutativeRingHom{f}{R,\oplus,\odot,0_R,1_R}{S,\boxplus,\boxtimes,0_S,1_S}}
}

\FormulaAxiomDeltaK[Quellstruktur eines Homomorphismus kommutativer Ringe]{%
  \CommutativeRingHom{f}{R,\oplus,\odot,0_R,1_R}
    {S,\boxplus,\boxtimes,0_S,1_S}
  \vdash \CommutativeRingAlg{R}{\oplus}{\odot}{0_R}{1_R}%
}{CommutativeRingHomSourceAxiom}{%
  \DeltaRow{Mengen}{R\dsep S\dsep0_R\dsep1_R\dsep0_S\dsep1_S}
  \DeltaRow{Funktionen}{f\dsep\oplus\dsep\odot\dsep\boxplus\dsep\boxtimes}
}

\FormulaAxiomDeltaK[Zielstruktur eines Homomorphismus kommutativer Ringe]{%
  \CommutativeRingHom{f}{R,\oplus,\odot,0_R,1_R}
    {S,\boxplus,\boxtimes,0_S,1_S}
  \vdash \CommutativeRingAlg{S}{\boxplus}{\boxtimes}{0_S}{1_S}%
}{CommutativeRingHomTargetAxiom}{%
  \DeltaRow{Mengen}{R\dsep S\dsep0_R\dsep1_R\dsep0_S\dsep1_S}
  \DeltaRow{Funktionen}{f\dsep\oplus\dsep\odot\dsep\boxplus\dsep\boxtimes}
}

\FormulaAxiomDeltaK[Ringhomomorphismusanteil eines Homomorphismus kommutativer Ringe]{%
  \CommutativeRingHom{f}{R,\oplus,\odot,0_R,1_R}
    {S,\boxplus,\boxtimes,0_S,1_S}
  \vdash
  \RingHom{f}{R,\oplus,\odot,0_R,1_R}
    {S,\boxplus,\boxtimes,0_S,1_S}%
}{CommutativeRingHomBaseHomAxiom}{%
  \DeltaRow{Mengen}{R\dsep S\dsep0_R\dsep1_R\dsep0_S\dsep1_S}
  \DeltaRow{Funktionen}{f\dsep\oplus\dsep\odot\dsep\boxplus\dsep\boxtimes}
}

\FormulaDefDeltaK[Isomorphismen kommutativer Ringe]{%
  \CommutativeRingIso{f}{R,\oplus,\odot,0_R,1_R}{S,\boxplus,\boxtimes,0_S,1_S}%
}{CommutativeRingIsoDef}{%
  \DeltaRow{Mengen}{R\dsep S\dsep0_R\dsep1_R\dsep0_S\dsep1_S}
  \DeltaRow{Funktionen}{f\dsep\oplus\dsep\odot\dsep\boxplus\dsep\boxtimes}
  \DeltaRow{\textbf{Axiome}}{}
  \DeltaPrem{Homomorphismus}{\CommutativeRingHom{f}{R,\oplus,\odot,0_R,1_R}{S,\boxplus,\boxtimes,0_S,1_S}}
  \DeltaPrem{Bijektivität}{f\colon R\bij S}
  \DeltaRow{\textbf{Neues Symbol}}{}
  \DeltaPrem{Isomorphismus}{\CommutativeRingIso{f}{R,\oplus,\odot,0_R,1_R}{S,\boxplus,\boxtimes,0_S,1_S}}
}

\FormulaAxiomDeltaK[Homomorphismusanteil eines Isomorphismus kommutativer Ringe]{%
  \CommutativeRingIso{f}{R,\oplus,\odot,0_R,1_R}
    {S,\boxplus,\boxtimes,0_S,1_S}
  \vdash
  \CommutativeRingHom{f}{R,\oplus,\odot,0_R,1_R}
    {S,\boxplus,\boxtimes,0_S,1_S}%
}{CommutativeRingIsoHomAxiom}{%
  \DeltaRow{Mengen}{R\dsep S\dsep0_R\dsep1_R\dsep0_S\dsep1_S}
  \DeltaRow{Funktionen}{f\dsep\oplus\dsep\odot\dsep\boxplus\dsep\boxtimes}
}

\FormulaAxiomDeltaK[Bijektivität eines Isomorphismus kommutativer Ringe]{%
  \CommutativeRingIso{f}{R,\oplus,\odot,0_R,1_R}
    {S,\boxplus,\boxtimes,0_S,1_S}
  \vdash f\colon R\bij S%
}{CommutativeRingIsoBijectiveAxiom}{%
  \DeltaRow{Mengen}{R\dsep S\dsep0_R\dsep1_R\dsep0_S\dsep1_S}
  \DeltaRow{Funktionen}{f\dsep\oplus\dsep\odot\dsep\boxplus\dsep\boxtimes}
}

Die Kommutativität der Multiplikation ist wiederum kein zusätzliches
Erhaltungsgesetz. Sie wird durch einen Ringisomorphismus in beide Richtungen
transportiert.

\chapter{Der kommutative Ring der ganzen Zahlen}

\FormulaThmDeltaK[Kommutativer Ring der ganzen Zahlen]{%
  \CommutativeRingAlg{\mathbb Z}{+}{\cdot}
    {\IntZero}{\IntOne}%
}{IntegerCommutativeRing}{%
  \DeltaRow{Mengen}{%
    \mathbb Z\dsep\IntZero\dsep\IntOne
    \dsep z\dsep w\dsep u}
  \DeltaRow{Funktionen}{+\dsep\cdot}
}
\begingroup
\ProofWideReasonsNearFormula
\begin{tabproofwide}
  \proofstepwidestarbreakkeep[]{\RingAlg{\mathbb Z}{+}{\cdot}{\IntZero}{\IntOne}}{%
    \FormulaRefAuto{IntegerUnitalRing}[thm]}
  \proofstepwidestarbreakkeep[2]{z\in\mathbb Z}{%
    \rA}
  \proofstepwidestarbreakkeep[3]{w\in\mathbb Z}{%
    \rA}
  \proofstepwidestarbreakkeep[2,3]{z\cdot w=w\cdot z}{%
    \FormulaRefAuto{IntegerMulCommutativeAxiom}[ax]{2,3}}
  \proofstepwidestarbreakkeep[]{\forall z\in\mathbb Z\,\forall w\in\mathbb Z\,(z\cdot w=w\cdot z)}{%
    \rUI{\rRI{2,\rUI{\rRI{3,4}}}}}
  \proofstepwidestarbreakkeep[]{\CommutativeRingAlg{\mathbb Z}{+}{\cdot}{\IntZero}{\IntOne}}{%
    \FormulaRefAuto{CommutativeRingDef}[def]{1,5}}
\end{tabproofwide}
\endgroup

\end{document}

